% sample-settings-viewport.tex -- one panel written to the standard in STYLE.md.
%
% Requires aethervk-manual.sty. Images live in img/ (written by the ManualCapture fixture, see
% PLAN.md). Until they exist, \uiscreenshot draws a placeholder of the given size, and the
% \uidefinecallout lines below place the markers at approximate positions. When the fixture
% has run, img/<figure>.callouts.tex replaces them with measured positions.
%
% Sources of every fact on this page (read 2026-11-01):
%   labels, layout, visibility  aethervk.ui-app/Views/SettingsTabView.axaml, lines 790-1210
%   help texts                  aethervk.ui-logic/Resources/AppStrings.resx (Tabs_Settings_*)
%   ranges and defaults         ViewportSettingsViewModel.cs, CameraService.cs
%                               (ViewportDisplayState.Default, Min/MaxNucleusExposureEv)
%   what is stored where        SaveFile/SaveIndexV1.cs (ViewportDto), PreferencesService.cs
%
% Items marked  % VERIFY  should be checked against the first capture.

% ---------------------------------------------------------------- callout positions
% One block per figure. The fixture writes exactly these commands, in the same order.
\InputIfFileExists{img/settings-viewport--orbit--light.callouts.tex}{}{%
  \uidefinecallout{settings-viewport}{CameraModeText}{1}{0.04}{0.07}%
  \uidefinecallout{settings-viewport}{DiagnosticPanelsCheckBox}{2}{0.04}{0.15}%
  \uidefinecallout{settings-viewport}{ComparisonLockCheckBox}{3}{0.04}{0.23}%
  \uidefinecallout{settings-viewport}{ProjectionPerspective}{4}{0.04}{0.31}%
  \uidefinecallout{settings-viewport}{ProjectionAdvancedExpander}{5}{0.04}{0.39}%
  \uidefinecallout{settings-viewport}{StarBackgroundBox}{6}{0.04}{0.48}%
  \uidefinecallout{settings-viewport}{OrbitLinesCheckBox}{7}{0.04}{0.56}%
  \uidefinecallout{settings-viewport}{ViewNucleusFromEarthButton}{8}{0.04}{0.65}%
  \uidefinecallout{settings-viewport}{SnapAbove}{9}{0.04}{0.82}}
\InputIfFileExists{img/settings-viewport-persp--orbit--light.callouts.tex}{}{%
  \uidefinecallout{settings-viewport-persp}{FovSlider}{1}{0.04}{0.12}%
  \uidefinecallout{settings-viewport-persp}{FovUnitPicker}{2}{0.04}{0.30}%
  \uidefinecallout{settings-viewport-persp}{PerspNearHelp}{3}{0.04}{0.48}%
  \uidefinecallout{settings-viewport-persp}{PerspFarHelp}{4}{0.04}{0.66}%
  \uidefinecallout{settings-viewport-persp}{DistanceUnitsHelp}{5}{0.04}{0.84}}
\InputIfFileExists{img/settings-viewport-ortho--orbit--light.callouts.tex}{}{%
  \uidefinecallout{settings-viewport-ortho}{OrthoHalfWidthSlider}{1}{0.04}{0.08}%
  \uidefinecallout{settings-viewport-ortho}{OrthoHalfHeightSlider}{2}{0.04}{0.19}%
  \uidefinecallout{settings-viewport-ortho}{OrthoUnitPicker}{3}{0.04}{0.30}%
  \uidefinecallout{settings-viewport-ortho}{OrthoNearHelp}{4}{0.04}{0.41}%
  \uidefinecallout{settings-viewport-ortho}{OrthoFarHelp}{5}{0.04}{0.52}%
  \uidefinecallout{settings-viewport-ortho}{OrthoProportionsLock}{6}{0.04}{0.66}%
  \uidefinecallout{settings-viewport-ortho}{ShowBehindCameraCheckBox}{7}{0.04}{0.77}%
  \uidefinecallout{settings-viewport-ortho}{RestoreOrthoProportionsButton}{8}{0.04}{0.90}}
\InputIfFileExists{img/settings-viewport-nucleus--nucleus--light.callouts.tex}{}{%
  \uidefinecallout{settings-viewport-nucleus}{PhysicalLightingToggle}{1}{0.04}{0.30}%
  \uidefinecallout{settings-viewport-nucleus}{NucleusExposureSlider}{2}{0.04}{0.80}}
\InputIfFileExists{img/settings-viewport-observer--model--light.callouts.tex}{}{%
  \uidefinecallout{settings-viewport-observer}{ObserverTopocentricCheckBox}{1}{0.04}{0.10}%
  \uidefinecallout{settings-viewport-observer}{ObserverNorthUpCheckBox}{2}{0.04}{0.19}%
  \uidefinecallout{settings-viewport-observer}{ObserverLatSlider}{3}{0.04}{0.29}%
  \uidefinecallout{settings-viewport-observer}{ObserverLonSlider}{4}{0.04}{0.38}%
  \uidefinecallout{settings-viewport-observer}{OrientFree}{5}{0.04}{0.50}%
  \uidefinecallout{settings-viewport-observer}{RestorePresetButton}{6}{0.04}{0.72}%
  \uidefinecallout{settings-viewport-observer}{EarthObserverScaleText}{7}{0.04}{0.86}}

% ================================================================= the panel
\begin{uipanel}[key=settings-viewport,
    open={\menu{Settings > Viewport 1} (one section per 3D viewport; the default layout has one)},
    modes={all three views; parts of the panel change with the view (Table~\ref{tab:viewport-modes})}]
  {Viewport}

The \uilabel{Viewport} section holds the settings of one 3D viewport's camera: how the scene is
projected onto the screen, what is drawn behind the comet, and, in the \uilabel{Model} view, where
on the Earth the observer stands and what the view tracks. None of these settings changes the
simulation. They change what the camera shows and how, and the dust already computed is
reused.

The section is named after its viewport (\uilabel{Viewport 1}, \uilabel{Viewport 2}, \dots).
Each 3D viewport added to the layout gets its own section with independent values. The
\uilabel{Show diagnostic panels} switch is the only exception: it applies to all viewports.

\paragraph{Opening the panel.} Select the \uilabel{Settings} tab and expand
\menu{Viewport 1}. It is the last section of the tab. The view-mode strip under the menu bar
(\uilabel{Orbit}, \uilabel{Nucleus}, \uilabel{Model}, \panelref{main}) and the view buttons on
the viewport overlay change the camera mode. The first line of this panel shows the current
mode.

\begin{table}[htbp]
  \centering\small
  \caption{Which parts of \menu{Settings > Viewport} appear in each view.
    \emph{Mode} is the text the panel prints in its first line.}
  \label{tab:viewport-modes}
  \begin{tabular}{@{}llccc@{}}
    \toprule
    Part of the panel & & \uilabel{Orbit} & \uilabel{Nucleus} & \uilabel{Model} \\
    \midrule
    \emph{Mode} line                           & & \uival{UpZenith} & \uival{CometOrbiting} & \uival{EarthPosition} \\
    Locks, \uilabel{Type}, \uilabel{Advanced: projection} & & \checkmark & \checkmark & \checkmark \\
    \uilabel{Background stars}, \uilabel{Orbit lines}, \uilabel{Nucleus from Earth} & & \checkmark & \checkmark & \checkmark \\
    \uilabel{Physical lighting}                & &            & \checkmark &            \\
    \uilabel{Nucleus exposure}                 & &            & \checkmark & \checkmark \\
    \uilabel{Snap above}                       & & \checkmark &            &            \\
    \uilabel{Earth Observer Settings}          & &            &            & \checkmark \\
    \bottomrule
  \end{tabular}
\end{table}
% (\toprule etc. need booktabs; the thesis preamble loads it -- VERIFY on the other computer)

\begin{figure}[htbp]
  \centering
  \uiscreenshot[placeholder={0.45,0.62}]{settings-viewport}{img/settings-viewport--orbit--light.png}
  \caption{\menu{Settings > Viewport 1} in the \uilabel{Orbit} view, with
    \uilabel{Advanced: projection} collapsed. Save \texttt{review\_67p\_2021-12}, light theme.}
  \label{fig:settings-viewport}
\end{figure}

% ----------------------------------------------------------------
\uisection[locks]{Mode and locks}{ViewportSettingsSection}

\begin{uiproperty}[key=mode, id=CameraModeText, type={read-only text},
    visible={always}]{Mode}
The camera mode of this viewport, printed as \uival{Mode: UpZenith},
\uival{Mode: CometOrbiting} or \uival{Mode: EarthPosition}. These are the internal names of
the \uilabel{Orbit}, \uilabel{Nucleus} and \uilabel{Model} views
(Table~\ref{tab:viewport-modes}). The rest of the panel adapts to this mode.
\end{uiproperty}

\begin{uiproperty}[key=diagnostics, id=DiagnosticPanelsCheckBox, type={check box},
    default={off}, effect=display, saved=prefs, help=Tabs_Settings_HelpViewDiagnostics]
  {Show diagnostic panels (all viewports)}
Shows the camera coordinates, the dust history, the observer line and the orbit debug panels
over every viewport at once. These panels are for checking the engine. A normal session
leaves them hidden.
\begin{uihelptext}
Shows the camera coordinates, dust history, observer line and orbit debug panels over the
viewport, in all viewports at once. For checking the engine; hidden by default.
\end{uihelptext}
\end{uiproperty}

\begin{uiproperty}[key=comparison-lock, id=ComparisonLockCheckBox, type={check box},
    default={off}, effect=camera, saved=none, help=Tabs_Settings_HelpViewComparisonLock]
  {Lock viewport for comparison}
Ignores mouse orbit, pan and zoom, and greys out \uilabel{Type} and
\uilabel{Advanced: projection}. The view then stays matched to a telescope image. The
billboard's \uilabel{Set up comparison} button (\menu{Settings > Billboards > FITS image})
turns the lock on as its last step. Clear the check box to move the camera again. The lock is
not saved: a reopened save starts unlocked.
\begin{uiexample}
After \uilabel{Set up comparison} the status line reads \uival{Comparison set up; the viewport
is locked.} and this box is checked. Clear it, orbit the camera to look at the tail from
the side, then use the billboard's \uilabel{Match camera to image} to return to the matched
view.
\end{uiexample}
\end{uiproperty}

% ----------------------------------------------------------------
\uisection[projection]{Projection}{}

\begin{uiproperty}[key=type, id=ProjectionPerspective, type={radio buttons},
    range={\uilabel{Perspective}, \uilabel{Orthographic}}, default={\uilabel{Perspective}},
    effect=camera, saved=file, visible={always; disabled while the viewport is locked},
    help=Tabs_Settings_HelpViewProjectionType]{Type}
\begin{description}
  \item[\uilabel{Perspective}] Distant objects look smaller, as through a camera lens. The
    field of view sets the zoom.
  \item[\uilabel{Orthographic}] No perspective: an object's size does not depend on its
    distance, as on a map or a flat telescope image. The half-width and half-height set what
    is visible.
\end{description}
The parameters of both are in \uilabel{Advanced: projection}. The second radio button has
the AutomationId \uiid{ProjectionOrthographic}.
\end{uiproperty}

\begin{uiproperty}[key=advanced, id=ProjectionAdvancedExpander, type={expander},
    default={collapsed}, help=Tabs_Settings_HelpViewProjectionAdvanced]{Advanced: projection}
Holds the field of view (perspective) or the half-width and half-height (orthographic), the
near and far clip planes and their display units. You rarely set these by hand. A billboard's
\uilabel{Match camera to image} sets the field from the image instead.
Figure~\ref{fig:settings-viewport-projection} shows both variants.
\end{uiproperty}

\begin{figure}[htbp]
  \centering
  \begin{minipage}[b]{0.48\linewidth}\centering
    \uiscreenshot[placeholder={0.95,0.62}]{settings-viewport-persp}{img/settings-viewport-persp--orbit--light.png}\\
    \small (a) \uilabel{Perspective}
  \end{minipage}\hfill
  \begin{minipage}[b]{0.48\linewidth}\centering
    \uiscreenshot[placeholder={0.95,1.05}]{settings-viewport-ortho}{img/settings-viewport-ortho--orbit--light.png}\\
    \small (b) \uilabel{Orthographic}
  \end{minipage}
  \caption{\menu{Settings > Viewport 1 > Advanced: projection} for each projection type.}
  \label{fig:settings-viewport-projection}
\end{figure}

\paragraph{Perspective parameters} (shown when \uilabel{Type} is \uilabel{Perspective}).
\uiusefigure{settings-viewport-persp}

\begin{uiproperty}[key=fov, id=FovSlider, type={slider (drag, or click and type)},
    unit={the one chosen in \uilabel{Units}: \unit{\degree}, \unit{\arcminute} or
      \unit{\arcsecond}},
    range={\qty{0.02}{\arcsecond} to \qty{179}{\degree}; step \qty{1}{\degree},
      \qty{1}{\arcminute} or \qty{10}{\arcsecond}},
    default={\qty{30}{\degree}}, % VERIFY: the engine's first projection may replace it
    effect=camera, saved=file, visible={Perspective; disabled while locked},
    help=Tabs_Settings_HelpViewFov]{FOV}
The \emph{vertical} field of view of the perspective camera. Smaller values zoom in.
Telescope fields are usually given in arcminutes, so switch \uilabel{Units} to
\uilabel{arcmin} before typing one.
\begin{uiexample}
A \qty{2048}{\pixel}-tall CCD frame at a plate scale of \qty{1.0}{\arcsecond} per pixel
spans \qty{2048}{\arcsecond}, about \qty{34.1}{\arcminute}. Set \uilabel{Units} to
\uilabel{arcmin}, click the slider and type \uival{34,1}, then press
\uikey{Enter}. The box uses the system's decimal separator, a comma on an Italian system.
\end{uiexample}
\end{uiproperty}

\begin{uiproperty}[key=fov-units, id=FovUnitPicker, type={drop-down list},
    range={\uilabel{deg}, \uilabel{arcmin}, \uilabel{arcsec}}, default={\uilabel{deg}},
    effect=display, saved=file, visible={Perspective},
    help=Tabs_Settings_HelpViewFovUnits]{Units}
The unit in which \uilabel{FOV} is shown and typed. It changes only the display: the field
itself stays the same. One arcminute is $1/60$ of a degree and one arcsecond is
$1/3600$ of a degree.
\end{uiproperty}

\begin{uiproperty}[key=persp-near, id=PerspNearHelp, type={slider, logarithmic},
    unit={the one chosen in \uilabel{Distance units}},
    range={\qty{1e-12}{\AU} to \qty{1000}{\AU}}, default={\qty{0.0001}{\AU}}, % VERIFY
    effect=camera, saved=file, visible={Perspective},
    help=Tabs_Settings_HelpViewPerspNear]{Near}
Objects closer to the camera than this distance are not drawn. The slider itself has no
AutomationId. The callout points at its \texttt{?} button (\uiid{PerspNearHelp}).
\end{uiproperty}

\begin{uiproperty}[key=persp-far, id=PerspFarHelp, type={slider, logarithmic},
    unit={the one chosen in \uilabel{Distance units}},
    range={\qty{1e-12}{\AU} to \qty{1000}{\AU}}, default={\qty{200}{\AU}}, % VERIFY
    effect=camera, saved=file, visible={Perspective},
    help=Tabs_Settings_HelpViewPerspFar]{Far}
Objects farther from the camera than this distance are not drawn.
\end{uiproperty}

\begin{uiproperty}[key=distance-units, id=DistanceUnitsHelp, type={drop-down list},
    range={\uilabel{Astronomical Units (AU)}, \uilabel{Kilometers (km)}, \uilabel{Meters (m)}},
    default={\uilabel{Astronomical Units (AU)}}, effect=display, saved=file,
    visible={Perspective}, help=Tabs_Settings_HelpViewDistanceUnits]{Distance units}
The unit used to show \uilabel{Near} and \uilabel{Far}. The values are stored in AU. This is
the same setting as \uilabel{Units} of the orthographic parameters: changing one changes the
other.
\end{uiproperty}

\paragraph{Orthographic parameters} (shown when \uilabel{Type} is \uilabel{Orthographic}).
\uiusefigure{settings-viewport-ortho}

\begin{uiproperty}[key=half-width, id=OrthoHalfWidthSlider, type={slider, logarithmic},
    unit={the one chosen in \uilabel{Units}},
    range={\qty{1e-12}{\AU} to \qty{1000}{\AU}}, default={\qty{0.0155}{\AU}}, % VERIFY
    effect=camera, saved=file, visible={Orthographic; disabled while locked},
    help=Tabs_Settings_HelpViewHalfWidth]{Half-Width}
Half the visible width of the view: it spans $\pm$ this value around its centre. With
\uilabel{Keep Proportions Locked} on, \uilabel{Half-Height} follows it.
\end{uiproperty}

\begin{uiproperty}[key=half-height, id=OrthoHalfHeightSlider, type={slider, logarithmic},
    unit={the one chosen in \uilabel{Units}},
    range={\qty{1e-12}{\AU} to \qty{1000}{\AU}}, default={\qty{0.0155}{\AU}}, % VERIFY
    effect=camera, saved=file, visible={Orthographic; disabled while locked},
    help=Tabs_Settings_HelpViewHalfHeight]{Half-Height}
Half the visible height of the view: it spans $\pm$ this value around its centre.
\begin{uiexample}
To frame a coma \qty{100000}{\kilo\metre} across, set \uilabel{Units} to
\uilabel{Kilometers (km)} and type \uival{50000} into \uilabel{Half-Height}. With the
proportions locked, \uilabel{Half-Width} follows from the viewport's shape.
\end{uiexample}
\end{uiproperty}

\begin{uiproperty}[key=ortho-units, id=OrthoUnitPicker, type={drop-down list},
    range={\uilabel{Astronomical Units (AU)}, \uilabel{Kilometers (km)}, \uilabel{Meters (m)}},
    default={\uilabel{Astronomical Units (AU)}}, effect=display, saved=file,
    visible={Orthographic}, help=Tabs_Settings_HelpViewOrthoUnits]{Units}
The unit for \uilabel{Half-Width}, \uilabel{Half-Height}, \uilabel{Near} and \uilabel{Far}.
It changes only the display. The perspective \uilabel{Distance units} list shows the same
setting.
\end{uiproperty}

\begin{uiproperty}[key=ortho-near, id=OrthoNearHelp, type={slider, logarithmic},
    range={\qty{1e-12}{\AU} to \qty{1000}{\AU}}, effect=camera, saved=file,
    visible={Orthographic}, help=Tabs_Settings_HelpViewOrthoNear]{Near}
The nearest distance along the view direction that is drawn. Anything closer is not drawn.
\end{uiproperty}

\begin{uiproperty}[key=ortho-far, id=OrthoFarHelp, type={slider, logarithmic},
    range={\qty{1e-12}{\AU} to \qty{1000}{\AU}}, effect=camera, saved=file,
    visible={Orthographic}, help=Tabs_Settings_HelpViewOrthoFar]{Far}
The farthest distance along the view direction that is drawn. Anything farther is not drawn.
\end{uiproperty}

\begin{uiproperty}[key=proportions-lock, id=OrthoProportionsLock, type={check box},
    default={on}, effect=camera, saved=file, visible={Orthographic},
    help=Tabs_Settings_HelpViewProportionsLock]{Keep Proportions Locked}
Keeps the ratio of half-width to half-height equal to the viewport's shape, also when the
window is resized, so the image is never stretched.
\end{uiproperty}

\begin{uiproperty}[key=behind-camera, id=ShowBehindCameraCheckBox, type={check box},
    default={on}, effect=display, saved=file, visible={Orthographic},
    help=Tabs_Settings_HelpViewShowBehindCamera]{Show objects behind camera}
Also draws the geometry behind the camera plane, within the \uilabel{Far} distance, as a
telescope image does. Dust behind the camera is never drawn, and a perspective view cannot
show it.
\end{uiproperty}

\begin{uiproperty}[key=restore-proportions, id=RestoreOrthoProportionsButton, type={button},
    effect=action, visible={Orthographic}]{Restore Proportions}
Recomputes \uilabel{Half-Width} from \uilabel{Half-Height} once, so the proportions match
the viewport again. Use it after a manual edit with the lock off. The text comes from the
button's tooltip. It has no \texttt{?} help.
\end{uiproperty}

% ----------------------------------------------------------------
\uiusefigure{settings-viewport}
\uisection[display]{Display}{}

\begin{figure}[htbp]
  \centering
  \uiscreenshot[placeholder={0.45,0.22}]{settings-viewport-nucleus}{img/settings-viewport-nucleus--nucleus--light.png}
  \caption{The parts of \menu{Settings > Viewport 1} that appear only in the
    \uilabel{Nucleus} view.}
  \label{fig:settings-viewport-nucleus}
\end{figure}

\uiusefigure{settings-viewport-nucleus}

\begin{uiproperty}[key=physical-lighting, id=PhysicalLightingToggle, type={check box},
    default={off}, effect=display, saved=prefs, visible={\uilabel{Nucleus} view},
    help=Tabs_Settings_HelpViewPhysicalLighting]{Physical lighting}
Lights the nucleus with direct sunlight only, scaled with the comet's heliocentric distance
$r$ as $(\qty{1}{\AU}/r)^2$, with no ambient fill. The exposure is set so that a white
surface facing the Sun at \qty{1}{\AU} is drawn white. The surface the Sun does not reach is
drawn in dark indigo. The choice is remembered across sessions
(Figure~\ref{fig:settings-viewport-nucleus}).
\end{uiproperty}

\begin{uiproperty}[key=nucleus-exposure, id=NucleusExposureSlider, type={slider},
    unit={\unit{\EV} (stops)}, range={\qty{-4}{\EV} to \qty{+6}{\EV}, step \qty{0.5}{\EV}},
    default={\qty{0}{\EV}}, effect=display, saved=none, % VERIFY: not in ViewportDto
    visible={\uilabel{Nucleus} and \uilabel{Model} views},
    help=Tabs_Settings_HelpViewNucleusExposure]{Nucleus exposure}
The brightness of the sunlit nucleus in photographic stops. Each whole step doubles or halves
it. The label above the slider shows the value, e.g.\ \uival{Nucleus exposure: +1.5 EV}. At
\qty{0}{\EV} a white surface facing the Sun is drawn white.
\begin{uiexample}
A nucleus at \qty{3}{\AU} receives $1/9$ of the sunlight it gets at \qty{1}{\AU}, about
\qty{3.2}{\EV} less. With \uilabel{Physical lighting} on, set the exposure to
\qty{+3}{\EV} to see its surface at about the brightness it would have at \qty{1}{\AU}.
\end{uiexample}
\end{uiproperty}

\uiusefigure{settings-viewport}% back to Figure~\ref{fig:settings-viewport}
\begin{uiproperty}[key=stars, id=StarBackgroundBox, type={drop-down list},
    range={\uilabel{Auto (off in Model view)}, \uilabel{On}, \uilabel{Off}},
    default={\uilabel{Auto (off in Model view)}}, effect=display, saved=prefs,
    help=Tabs_Settings_HelpViewStarBackground]{Background stars}
A decorative procedural star field. It does not match the real sky, and its field of view
is fixed. \uilabel{Auto} draws it in every view except \uilabel{Model}, which then has a
black backdrop. \uilabel{On} and \uilabel{Off} apply to
every view.
\end{uiproperty}

\begin{uiproperty}[key=orbit-lines, id=OrbitLinesCheckBox, type={check box}, default={on},
    effect=display, saved=prefs, help=Tabs_Settings_HelpViewOrbitLines]{Orbit lines}
Draws the comet's reference orbit (red), the Earth's orbit (light blue) and the path
travelled in this run (yellow). In the \uilabel{Orbit} and \uilabel{Model} views the overlay
legend lists them (\panelref{overlay}).
\end{uiproperty}

\begin{uiproperty}[key=nucleus-from-earth, id=ViewNucleusFromEarthButton, type={button},
    effect=action, visible={always; enabled once a comet is committed}, % VERIFY enable rule
    help=Tabs_Settings_HelpViewNucleusFromEarth]{Nucleus from Earth}
Switches to the Earth observer tracking the comet, with celestial north up and a field a few
nucleus diameters tall. The result is the sunlit nucleus as seen from the Earth. It is a
shortcut for: \uilabel{Model} view, \uilabel{Comet tracking}, \uilabel{North up} on, then a
small \uilabel{FOV}.
\end{uiproperty}

\begin{uiproperty}[key=snap-above, id=SnapAbove, type={three buttons},
    range={\uilabel{Sun}, \uilabel{Comet}, \uilabel{Earth}}, effect=action,
    visible={\uilabel{Orbit} view}]{Snap above}
Flies the camera above the chosen body and fits the projection to it. \uilabel{Sun} restores
the startup view. \uilabel{Comet} frames the committed comet's nucleus. \uilabel{Earth} fits
the Earth's radius. This is an animation only: the camera does not track the body
afterwards. The buttons have no AutomationId. The fixture finds them by name inside
\uiid{ViewportSettingsSection} (synthetic id \uiid{SnapAbove}).
\end{uiproperty}

% ----------------------------------------------------------------
\uiusefigure{settings-viewport-observer}
\uisection[observer]{Earth Observer Settings}{}

Shown only in the \uilabel{Model} view, where the camera stands at the Earth. The settings
below are saved with the file, because a comparison with a telescope image depends on them.

\begin{figure}[htbp]
  \centering
  \uiscreenshot[placeholder={0.45,0.70}]{settings-viewport-observer}{img/settings-viewport-observer--model--light.png}
  \caption{\uilabel{Earth Observer Settings} in the \uilabel{Model} view, observer on the
    surface, \uilabel{Comet tracking}. The save's site is used.}
  \label{fig:settings-viewport-observer}
\end{figure}

\begin{uiproperty}[key=topocentric, id=ObserverTopocentricCheckBox, type={check box},
    default={off (geocentric)}, effect=camera, saved=file,
    help=Tabs_Settings_HelpObserverTopocentric]{Observer on the Earth's surface (topocentric)}
\begin{description}
  \item[On (topocentric)] The observer stands at the telescope site given by
    \uilabel{Latitude (°N)} and \uilabel{Longitude (°E)} and turns with the Earth. Nearby
    objects shift by diurnal parallax, and a horizon appears.
  \item[Off (geocentric)] The camera is at the Earth's centre and follows only its motion
    around the Sun.
\end{description}
The view never rotates with the Earth.
\end{uiproperty}

\begin{uiproperty}[key=north-up, id=ObserverNorthUpCheckBox, type={check box}, default={on},
    effect=camera, saved=file, help=Tabs_Settings_HelpObserverNorthUp]{North up}
On: the view is levelled to celestial north (J2000), with north up and east left, like a
telescope image. Off: the view is levelled to the ecliptic pole on the site's side of the
equator.
\end{uiproperty}

\begin{uiproperty}[key=latitude, id=ObserverLatSlider, type={slider}, unit={\unit{\degree}, positive north},
    range={\qty{-90}{\degree} to \qty{90}{\degree}, step \qty{0.1}{\degree}},
    default={\qty{0}{\degree}}, effect=camera, saved=file, visible={topocentric on},
    help=Tabs_Settings_HelpObserverLatitude]{Latitude (°N)}
The latitude of the observing site.
\end{uiproperty}

\begin{uiproperty}[key=longitude, id=ObserverLonSlider, type={slider}, unit={\unit{\degree}, positive east of Greenwich},
    range={\qty{-180}{\degree} to \qty{180}{\degree}, step \qty{0.1}{\degree}},
    default={\qty{0}{\degree}}, effect=camera, saved=file, visible={topocentric on},
    help=Tabs_Settings_HelpObserverLongitude]{Longitude (°E)}
The longitude of the observing site.
\begin{uiexample}
For an observatory at \ang{45;3;}\,N, \ang{7;40;}\,E, type \uival{45,05} into
\uilabel{Latitude (°N)} and \uival{7,67} into \uilabel{Longitude (°E)}. Sites west of
Greenwich take negative longitudes.
\end{uiexample}
\end{uiproperty}

\begin{uiproperty}[key=look-direction, id=OrientFree, type={radio buttons},
    range={\uilabel{Free look}, \uilabel{Comet tracking}, \uilabel{Sun tracking}},
    default={\uilabel{Free look}}, effect=camera, saved=file,
    help=Tabs_Settings_HelpObserverLookDirection]{Look Direction}
\begin{description}
  \item[\uilabel{Free look}] Drag to look around. The direction stays fixed among the stars.
  \item[\uilabel{Comet tracking}] Keeps the comet exactly at the centre of the view
    (\uiid{OrientCometTracking}).
  \item[\uilabel{Sun tracking}] Keeps the Sun at the centre of the view
    (\uiid{OrientSunTracking}).
\end{description}
On the surface, while the tracked target is below the horizon, the view faces the point on
the horizon where it will rise. A box then explains the hold and offers
\uilabel{Move observer there}. That button moves the site to the latitude and longitude
where the target is overhead now.
\end{uiproperty}

\begin{uiproperty}[key=restore-preset, id=RestorePresetButton, type={button}, effect=action,
    visible={enabled after the projection was changed away from the tracking preset}]
  {Restore preset projection}
Re-applies the tracking projection, in which the target spans \qty{10}{\percent} of the
view.
\end{uiproperty}

\begin{uiproperty}[key=scale, id=EarthObserverScaleText, type={read-only text},
    visible={Model view, once a scale is known}, help=Tabs_Settings_HelpObserverScale]
  {Telescope scale}
The angle one screen pixel covers and the field of view. At the comet's distance it also
gives the kilometres per pixel and the half-width and half-height of the view, for example:
\begin{center}\small
\uival{Scale 1.2″/px · field 21.26′ × 12′ (1063 × 600 px)}\\
\uival{At the comet (0.4664 AU): 406 km/px, half-width 2.1579E+05 km, half-height 1.218E+05 km}
\end{center}
The figures above are illustrative. The text always uses a decimal point, whatever the
system's language.
\end{uiproperty}

\paragraph{See also.} \menu{Settings > Billboards > FITS image} for
\uilabel{Match camera to image} and \uilabel{Set up comparison}, which set many of the values
above in one step. The \uilabel{Model} button of the view-mode strip (\panelref{main}).

\end{uipanel}
